\thesistitle{A High and Variable Dimensional Measurement of the $Z$+jets Differential Cross Section with the ATLAS Experiment and Artificial Intelligence}

%"Dissertation" for PhD, "Thesis" for master's
\documenttitle{Dissertation}

\degreename{Doctor of Philosophy}

% Use the wording given in the official list of degrees awarded by UCI:
% http://www.rgs.uci.edu/grad/academic/degrees_offered.htm
\degreefield{Physics and Astronomy}

% Your name as it appears on official UCI records.
\authorname{Kevin Greif}

% Use the full name of each committee member and full title 
% (e.g. Professor/Associate Professor).
\committeechair{Daniel Whiteson}
\othercommitteemembers
{
  Andrew Lankford \\
  Jonathan Feng
}

\degreeyear{2026}

\copyrightdeclaration
{
  {\copyright} {\Degreeyear} \Authorname
}

% If you have previously published parts of your manuscript, you must list the
% copyright holders; see Section 3.2 of the UCI Thesis and Dissertation Manual.
% Otherwise, this section may be omitted.
% \prepublishedcopyrightdeclaration
% {
% 	Chapter 4 {\copyright} 2003 Springer-Verlag \\
% 	Portion of Chapter 5 {\copyright} 1999 John Wiley \& Sons, Inc. \\
% 	All other materials {\copyright} {\Degreeyear} \Authorname
% }

% The dedication page is optional
% (comment out to exclude).
\dedications
{  
  To my parents, Douglas and Anne Greif
}

\acknowledgments
{
  I have many, many people to thank. I'll fill this out soon!
  
  % (You must acknowledge grants and other funding assistance. 
  
  % You may also acknowledge the contributions of professors and
  % friends.
  
  % You also need to acknowledge any publishers of your previous
  % work who have given you permission to incorporate that work
  % into your dissertation. See Section 3.2 of the UCI Thesis and
  % Dissertation Manual.)
}


% Curriculum vitae is in vita.tex, which is \input into main.tex.

% ============================================================
% Statement on the use of AI tools
%
% Rendered as a preliminary page right after the Vita and just
% before the Abstract. The injection point is configured in
% main.tex via a small patch to \curriculumvitaepage, so this
% page appears inside the standard \preliminarypages flow and
% gets a TOC entry.
%
% Fill in the introductory paragraph and table rows below.
% Each AI-use cell should keep just one of {Low, Medium, High}
% (delete the others). Add or remove rows as appropriate.
% ============================================================
\newcommand{\aiusestatementpage}{%
  \begin{center}
    \textbf{\Large STATEMENT ON THE USE OF AI TOOLS}
  \end{center}
  \parskip 12pt
  \parindent 0pt

  AI tools were used frequently in the underlying research documented in this thesis, especially for writing and editing software.
  Additionally AI tools were used for the ideation and background research needed to compile this thesis, especially for Chapters 2 and 3.
  On the level of the text, AI tools were used for organizing and outlining various sections but all of the sentences appearing in the final version of the thesis were written by hand.
  Many of the diagrams in this thesis were built with AI tools, and are labeled as such in the figure captions.

  Some overviews of the particular AI-based workflows I used when drafting this thesis are provided in the table below.
  Models used were Claude Sonnet 4.6 and 5.0, Opus 4.6 and 4.7, and GPT 5.4 and 5.5.
  Agentic tools used were Claude Code, Claude Cowork, Cursor, and OpenClaw.

  \vspace{0.5em}
  \begin{singlespace}
  \noindent\begin{tabularx}{\textwidth}{@{}
      >{\raggedright\arraybackslash}p{0.24\textwidth}
      >{\centering\arraybackslash}p{0.13\textwidth}
      X@{}}
    \toprule
    \textbf{Task} & \textbf{AI use} & \textbf{Example workflow} \\
    \midrule
    Section outlining and organization
      & Light
      & Prompt AI to translate ideas and bullet points into \LaTeX{} documents for further editing \\
    \addlinespace
    Copy editing
      & High
      & Prompt AI to read a recently written section and fix any spelling and grammar issues \\
    \addlinespace
    Diagram design
      & High
      & Prompt AI to write code for producing diagrams with Matplotlib or TikZ \\
    \addlinespace
    Literature search
      & Low
      & Prompt AI to search for relevant citations from the InspireHEP database \\
    \addlinespace
    Concept review
      & Med
      & Prompt AI for explaining and clarifying QCD concepts \\
    \addlinespace
    \LaTeX{} formatting
      & High
      & Promt AI to debug compilation issues, add Figures to documents, format tables, etc. \\
    \bottomrule
  \end{tabularx}
  \end{singlespace}
  \clearpage
}

% The abstract was previously limited to a maximum of 350 words, 
% but the UCI manual at https://etd.lib.uci.edu/electronic/td2e#2.2.1.
% currently does not indicate that there is any word limit for the abstract
\thesisabstract
{
  Proton-proton collisions at the Large Hadron Collider (LHC) offer the opportunity to observe the interactions of fundamental particles at very high energy scales. The high collision rate and large number of final state particles produced per collision imply that the datasets produced by detectors such as the ATLAS experiment are large and highly dimensional. The complexity of these datasets calls for the use of novel data analysis techniques which can exploit all of the available information to illuminate known interactions and search for new ones. This thesis presents applications of artifical intelligence (AI) techniques to improve the physics results of the ATLAS experiment. It centers on a full phase measurement of the Z+jets production cross section at the LHC, where the cross section is measured differential in the kinematics of every final state charged particle through the use of an AI based unfolding algorithm. This is the first such measurement performed at the LHC, which provides a complete experimental characterisation of the Z+jets production process and a proof-of-principle for the further pursuit of such measurements on a host of LHC processes.
}


%%% Local Variables: ***
%%% mode: latex ***
%%% TeX-master: "thesis.tex" ***
%%% End: ***
